\chapter[Modeling BSS Fleet Mobility and Battery Status]{Modeling \ac{bss} Fleet Mobility and Battery Status}
\label{ch:mobility_and_battery_modeling_bicing}

Building on Chapter~\ref{ch:bicing_mobility_patterns}, which characterized the mobility patterns of \textit{Bicing}, this Chapter addresses a central operational challenge in electrified docked systems: forecasting battery status. Low state of charge removes \acp{e-b} from the available fleet and can intensify spatial imbalances. At the same time, operators rarely have continuous GPS routes or in-transit battery telemetry, so prediction must rely on the trip-based operational data that \acp{bss} routinely collect. Therefore, this Chapter addresses the following research question: \emph{How can \ac{e-b} battery levels be predicted without exact routes or in-transit measurements, so as to improve bike availability and operational efficiency?}

To address the research question posed above, this Chapter starts from the fact that an \ac{e-b}'s battery level depends on the trips it makes and on the time it spends docked between them, so forecasting battery status first requires forecasting where and when each bike will travel. It therefore develops a joint framework with two components: a Markov-chain mobility model that predicts the future locations of each bike, and a battery model that estimates how its charge evolves along those trips, combining physics-inspired consumption with linear charging at stations. The Chapter first reviews the related literature and documents the spatiotemporal variation in battery levels in \textit{Bicing}, which motivates this approach. It then presents both models and shows how their combination supports bike-level and system-level forecasts of locations and battery levels, informing charging and rebalancing decisions.

The material in this Chapter is adapted from Bassolas et al.~\cite{Bassolas2025}, published in \textit{Scientific Reports}.

\section{Introduction}

By 2022, more than 2,000 \acp{bss} were available worldwide, with approximately 9 million bikes in circulation~\cite{MeddinWorldMap2025}. These systems enhance the efficiency of \ac{pt} by improving access to stations~\cite{Noland2019, Kapuku2021}, which in turn reduces car usage and alleviates traffic congestion~\cite{Ma2020, Fan2020}. Moreover, the increased use of bikes and the resulting shift in transport modes have been shown to positively impact public health by promoting physical activity~\cite{Otero2018}, benefit the environment by reducing pollutant emissions~\cite{Clockston2021}, and lead to economic savings~\cite{Ricci2015}.

The recent increase in data availability related to \acp{bss} has led to greater interest in understanding the mobility patterns within these systems~\cite{Kou2019, Li2020}. Studies have shown that stations with strong flows tend to be located nearby~\cite{Chen2017}, and they often display a certain level of mesoscopic organization within communities~\cite{Zaltz2013}. Additionally, the development of models to understand the usage of micro-mobility~\cite{Tran2015} has improved understanding of factors such as topography~\cite{Kim2020}, day of the week~\cite{Tran2015}, weather conditions~\cite{ElAssi2017, Kim2018, Ashqar2019}, and the built environment~\cite{Tran2015}. Furthermore, efforts have been made to develop predictive models that enhance the management of \acp{bss} infrastructures efficiently and guide system expansion through \ac{ml} and probabilistic approaches~\cite{Singhvi2015, Li2015, Chen2016, Albuquerque2021, Yang2020graph, Almannaa2020}. In this line, one of the major challenges in optimizing \ac{bss} performance is ensuring bike availability, as stations often experience significant flow imbalances~\cite{Yang2016, Singla2015, FaghihImani2017Empirical}. To address this, models focused on optimizing rebalancing operations have been developed~\cite{Raviv2013, Dell2014, Dell2018, Chiariotti2018}.

In recent years, \acp{bss} have undergone progressive electrification~\cite{MeddinWorldMap2025, Galatoulas2020}, with a sharp increase in the number of \acp{e-b} compared to their mechanical counterparts~\cite{Galatoulas2020, Bielinski2020}, as \acp{e-b} allow for longer trips and attract a broader user base~\cite{Simsekoglu2019, Weinert2007}. \acp{e-b} have been shown to be less sensitive to factors such as longer distances, adverse weather conditions, and casualties~\cite{Wamburu2021, Campbell2016, Siman2018}. Despite these advantages, \acp{e-b} present additional challenges to maintaining bike availability, primarily due to low battery levels, battery malfunctions, and complications arising from high temperatures~\cite{Ji2014, Florez2018, Zhang2019, Usama2019}. Moreover, the scarcity of battery data at the system level complicates the development of predictive models. While several studies have implemented frameworks to predict the battery consumption of individual bikes during trips~\cite{Burani2022, Steyn2014}, most approaches at the \ac{bss} level focus on the station-level granularity~\cite{Zhou2023}.

This publication addresses the challenge that the state of charge poses to \ac{e-b} availability by analyzing mobility and battery level data from Barcelona's \ac{bss}, \textit{Bicing}, and proposing a model that predicts future bike locations with their corresponding battery levels. The structure of this Chapter is as follows: first, the spatial and temporal variability in average battery levels is analyzed to identify regions more vulnerable to battery depletion. Next, a Markov chain approach is implemented to predict the future locations of bikes. Subsequently, the parameters in the equations describing battery power consumption are optimized. Finally, by combining the models presented in the two previous Sections, a framework is provided to support the management of \ac{e-b} fleets by projecting future scenarios.
\section{Methodology}
\label{sec:battery_methodology}

    \subsection{Data}

    This research analyzes the bike usage and battery levels of Barcelona's \ac{bss} \cite{Soriguera2018, Winslow2019, Grau-Escolano2024}, using data supplied by its operator, \textit{Bicing}, which made available the following data sets:

      \begin{itemize}
        \item The trips dataset, described in detail in Chapter~\ref{ch:bicing_mobility_patterns}, is used in its entirety for this Chapter, covering the period from April 2019 to December 2022. The same trip filtering as in that Chapter is applied, retaining only trips with duration between 2 and 60 minutes. In this work, we focus specifically on the mobility of \acp{e-b}.

        \item Two \ac{e-b} battery dynamics datasets. The first one, with snapshots taken every 8 hours over the entire year of 2022, and the second one, with snapshots taken every 30 minutes over three weeks from 29 March to 18 April 2023. Each snapshot was recorded at the system level, capturing the battery levels of every bike available at the docking stations at the specified intervals. As the snapshots were captured at the stations, battery levels were recorded exclusively while the bikes were docked and not during travel. The relevant fields in both datasets include percentage, voltage, current, and temperature of the battery, bike station, timestamp, bike status (operative, low battery, or inoperative for several reasons such as maintenance), and location (at station, in use, or at the workshop). A detailed analysis of these datasets is provided in Appendices~\ref{app_paper2:data_description}.
      \end{itemize}

    Additionally, trip data from March and April 2023 was available to assess battery consumption in the 30-minute battery dataset.



    \subsection[Predicting e-bikes' locations]{Predicting \acp{e-b}' locations}

    Although the datasets record where each bike is at a given time, forecasting battery levels requires knowing where bikes will be in the future. To predict these future bike locations, a probabilistic model based on a Markov chain approach has been developed \cite{Lu2013, Huang2015}. In this model, each bicycle of the system is modeled independently, with transitions based on the probability of movement from its current station. Given a starting location for a bike and a forecasting horizon $\Delta t$, the algorithm follows the sequence below:

    \begin{itemize}
        \item An inter-event time is sampled from the distribution $I_i^{w,h}$ to simulate the resting time of a bike given a station $i$ on a specific day type $w$ and hour $h$. The days of the week are grouped into four categories: Working days (Monday to Thursday), Friday, Saturday, and Sunday. The distribution $I_i^{w,h}$ is calculated by determining the time elapsed between a bike’s arrival at station $i$ and its next departure.
        
        \item If the inter-event time exceeds the desired $\Delta t$, the process stops, and the final location of the bike will correspond to its current station. 
        
        \item If the forecasting horizon is greater than the inter-event time, a destination $j$ is sampled based on the probability $M_{ij}^{w,h}$, given a weekday $w$, hour $h$, and origin station $i$. This destination then becomes the new location of the bike, and the average travel time between the station pairs, $\tau_{ij}$, is elapsed. The destination probability is calculated using the number of trips $T_{ij}^{w,h}$ between $i$ and $j$ from the 2022 trip data, as follows:
        
            \begin{equation}
                M_{ij}^{w,h}=\frac{T_{ij}^{w,h}}{\sum_{\forall j} T_{ij}^{w,h}}, 
                \label{eq:destination_probability}
            \end{equation}

        \item If the total time elapsed after completing the trip exceeds the desired $\Delta t$, the final predicted location will be station $j$. Otherwise, the process starts again.
        
    \end{itemize}

    For each bike, 200 realizations are performed to capture a representation of all possible scenarios. The main result is the probability of ending up at station $j$ after a time period $\Delta t$, denoted as $f^{w,h}_{i,j}(\Delta t)$. This probability is calculated as the number of times a bike ends up at station $j$, divided by the total number of iterations. To aggregate the results at the system level and determine the number of bikes available per station, the probabilities $f^{w,h}_{i,j}(\Delta t)$ are summed for each station $j$, which may result in non-integer values for the number of bikes at each destination. The initial conditions required by the model are the station location of each bike in \ac{bss} together with the corresponding day of the week and initial hour.



    \subsection{Mathematical formulation of battery modeling}
    \label{sec:battery_modeling}

      The battery levels of the \textit{Bicing} \acp{e-b} are modeled with a method based on the fundamental laws of motion. The instantaneous power consumed while riding is derived from the need to overcome three main forces opposing the motion: gravity resistance ($F_g$), rolling resistance ($F_r$), and aerodynamic drag ($F_a$) \cite{Burani2022, Steyn2014}.

      The gravity resistance $F_g$ is given by:
      \begin{equation}
      F_g=g \cdot m \cdot \sin(arctan(s)),
      \end{equation}
      where $g$ is the gravitational acceleration, $m$ is the combined mass of the rider and bike, and $s$ represents the slope. Similarly to \cite{Burani2022}, $F_g=0$ when $F_g\leq 0$.

      Rolling resistance $F_r$ is given by:
      \begin{equation}
      F_r={C_{rr}} \cdot g \cdot m  \cdot \cos(arctan(s)),
      \end{equation}
      where $C_{rr}$ is the rolling resistance coefficient.

      Finally, the aerodynamic drag $F_a$ is given by:
      \begin{equation}
      F_a=0.5 \cdot C_d \cdot A \cdot \rho \cdot (v+w)^{2},
      \end{equation}
      where $C_d$ denotes the drag coefficient, $A$ is the frontal area, $\rho$ is the air density, and $v$ and $w$ are the bike and wind speeds, respectively. 

      The instantaneous power consumed is then given by:
      \begin{equation}
      P(t)=vF,
      \end{equation}
      Assuming no energy losses during the process, the total power consumed over a trajectory of duration $t$ is given by $tv(F_g+F_r+F_a)$. As the instantaneous speed information is unavailable, the speed is approximated by dividing the total distance traveled by the trip time. Consequently, the energy consumed in a trip is given by $d_{ij}(F_g+F_r+F_a)$. The distance traveled $d_{ij}$ is approximated using the OpenRouteService \ac{api} \cite{OpenRouteService_2023}, which recommends bicycle routes based on \ac{osm} data \cite{OpenStreetMap_2023}. The slope $s$ is derived from elevation data obtained from \ac{otd} \cite{OpenTopoData_2023}, with station altitude details provided in Appendix~\ref{app_paper2:subsection_elevation_data}. To convert power consumed into a battery percentage, and taking into account that the \acp{e-b} only provide support to the rider when pedaling, a constant $K_{\rm trip}$ has been added by multiplying the whole equation $K_{\rm trip}d_{ij}(F_g+F_r+F_a)$.

      On the other hand, the rate of battery recharge during rest periods is modeled as:
      \begin{equation}
      \Delta B_{\rm rest}=\frac{K_{\rm rest}t_{\rm rest}}{60}, \label{charge}
      \end{equation}
      where $K_{\rm rest}$ represents the charging constant, and $t_{\rm rest}$ is the rest time. The inverse of $K_{\rm rest}$ corresponds to the time required to fully charge a bike from $0\%$ to $100\%$. The division by $60$ is included to express $K_{\rm rest}$ in units of $min^{-1}$. All variables are represented in the International System of Units.

      In summary, the battery model employs a linear approach: battery levels increase linearly with time when bikes are docked (Eq.~\ref{charge}), and decrease during trips according to the consumption function described below (Eq.~\ref{consumption}). The main variables used to compute battery consumption during a trip are its duration, the distance between the origin and destination stations, and the slope along the route, with the consumption calculated by integrating the instantaneous power consumption over the trip duration.
    
      
      
    \subsection{Optimization of the battery dynamics model}
    \label{sec:battery_optimization}

      The battery consumption during a trip can be approximated by an expression that groups the relevant constants. This expression is written as:

        \begin{equation}
          \Delta B_{\rm trip}=K_{\rm trip} d_{ij}(K_a \sin(\arctan(s))+ K_b \cos(\arctan(s)) + K_c v^2), \label{consumption}
        \end{equation}
      where
        \begin{align*}
          & K_a=g \cdot m, \\
          & K_b={C_{rr}} \cdot g \cdot m, \\
          & K_c=0.5 \cdot C_d \cdot A \cdot \rho.
        \end{align*}

      Although the constant $K_{\rm trip}$ is not strictly required mathematically, it allows us to explore feasible values for $K_a$, $K_b$, $K_c$. Apart from $g$, the constants involved in the computation of $K_a$, $K_b$, and $K_c$ are unknown and may vary depending on factors such as rider characteristics, weather conditions, bike status, and terrain type. Therefore, an optimization process is necessary to determine the optimal values of $K_a$, $K_b$, and $K_c$ that enable accurate battery consumption predictions, even when the constants of every trip are not known precisely.

      To perform the optimization, a different range of values is considered for each constant. The mass $m$ is assumed to range between 79 kg and 119 kg, considering the bike's weight of 29 kg. Based on previous studies \cite{Steyn2014, Bertucci2013}, the rolling resistance coefficient can range between $0.001$ and $0.003$, the drag coefficient between $0.6$ and $0.8$, and the effective area between $0.4$ and $0.6$. The air density $\rho$ is fixed at $1.225kg/m^3$, corresponding to a temperature of $15^{\circ}C$ at sea level. Given these constants' value ranges, the following limits are defined for the hyperparameter optimization: $K_a = [774.2, 1166.2]$, $K_b = [0.7742, 3.4986]$, and $K_c = [0.147, 0.294]$. The range for $K_{\rm rest}$ is set between $[0.002083, 0.005556]$, corresponding to full-charging times from 480 to 180 minutes, respectively. The range for $K_{\rm trip}$ is set between $[1 \cdot 10^{-8},1]$.

      The Optuna framework~\cite{Optuna2019} was utilized for hyperparameter optimization, minimizing the root mean squared error (RMSE) between observed and predicted battery levels and using 75\% of the three-week 30-minute dataset for training and 25\% for validation.
\section{Results}
\label{sec:battery_results}

    \subsection[Assessing e-bikes' battery levels and usage patterns]{Assessing \acp{e-b}' battery levels and usage patterns}
    \label{subsec:battery_results_analysis}

        The average battery levels per neighborhood show significant variability across neighborhoods (Figure~\ref{battery_overiew}A, Appendix~\ref{app_paper2:subsection_battery_levels_across_time}). To facilitate the analysis of all the temporal series of average battery levels derived from each neighborhood, K-means clustering was performed, resulting in four distinct clusters (Figure~\ref{battery_overiew}B, Appendix~\ref{app_paper2:subsection_choice_of_number_of_clusters}). The central and southern neighborhoods, which fall into Clusters 1 and 4, have the lowest battery levels, exhibiting steep declines in battery levels during business hours with recovery during nighttime (Figure~\ref{battery_overiew}C). This is likely due to the high levels of industrial and commercial activity, as well as the dense populations in these areas. Additionally, it is worth noting that some neighborhoods in Cluster 4 have slower charging speeds at night, likely due to a stronger nightlife presence. In contrast, northern neighborhoods in Clusters 2 and 3 are mainly residential with less commercial activity, resulting in higher and more stable battery levels throughout the day.

        \begin{figure}[!htbp]
          \centering
          \includegraphics[width=0.9\textwidth]{0_plots/paper2/main/Panel1.png}
          \caption[Average neighborhood battery levels and temporal clustering]{
            \textbf{Average neighborhood battery levels and temporal clustering.}
            \textbf{(A)} Average battery levels across neighborhoods. The average for each neighborhood is computed using the battery values of all the bikes in that neighborhood. 
            \textbf{(B)} Clustering of the neighborhoods according to the temporal pattern of battery levels using K-means into four distinct clusters. 
            \textbf{(C)} Temporal series of average battery levels for each cluster. The neighborhoods in grey do not have bike-sharing stations.}
          \label{battery_overiew}
        \end{figure}

        One factor directly related to average battery levels is the average time between trips for each bike, which can also be referred to as inter-event time. With longer inter-event times, \acp{e-b} have more opportunity to recharge, resulting in higher battery levels. In Figure~\ref{correlations}A, central and southern neighborhoods (Clusters 1 and 4) have not only lower average battery levels but also shorter inter-event times, indicating a higher turnover of \acp{e-b} due to frequent trips. This high turnover aligns with the higher density of users and increased mobility demand in commercial and densely populated areas. However, Zona Franca, located in the southernmost area of the city, presents a distinct case due to its industrial nature and significantly low population density (1 inhabitant per hectare \cite{bcn_population_density}), which results in low \ac{e-b} usage. Conversely, northern neighborhoods (Clusters 2 and 3) exhibit longer inter-event times, reflecting less frequent \ac{e-b} usage and consequently higher and more stable battery levels despite some of them being in high altitudes. Additionally, a Pearson correlation analysis was conducted between average battery levels and inter-event time at the station level (Figure~\ref{correlations}B), revealing a correlation value of 0.39. This moderate positive correlation suggests that stations with longer inter-event times tend to have higher average battery levels, further supporting the neighborhood-level results.

        \begin{figure}[!htbp]
            \centering
            \includegraphics[width=0.9\textwidth]{0_plots/paper2/main/Panel2.png}
            \caption[Inter-event times, distance index, and their relationships with battery levels]{
              \textbf{Inter-event times, distance index, and their relationships with battery levels.} 
              \textbf{(A)} Average inter-event time between trips across neighborhoods. The average for each neighborhood is computed using the inter-event time values of all bikes, with outliers removed using the IQR method at the station level. 
              \textbf{(B)} Relation between average battery levels and inter-event time at the station level, with a Pearson correlation of 0.39 (p-value<0.001). 
              \textbf{(C)} Distance index per neighborhood, calculated by first computing the index at the station level to accurately reflect trip distances between stations, and then averaging it at the neighborhood level. 
              \textbf{(D)} Relation between average battery levels and the distance index at the station level, with a Pearson correlation of 0.47 (p-value<0.001). The neighborhoods in gray do not have bike-sharing stations.}
            \label{correlations}
        \end{figure}

        Apart from the inter-event time, two additional factors can affect the average battery levels at a station: the origins of the trips that the station receives and the number of trips coming from each origin. To evaluate the impact of these factors on battery depletion, a distance index (\(DI_j\)) was developed for each station \( j \). This index measures the influence of the distances traveled by bikes on their battery levels by accounting for both the distance of each trip and the frequency of trips from each origin station. Specifically, for trips originating at station \( i \) and ending at station \( j \), the distance  (\( d_{i,j} \)) is weighted by the inverse of the total trips originating from station \( i \) (\( n_i \)). This weighting ensures that trips from less frequent origins have a greater influence on the distance index. The weighted distances for all origin stations \( i \) are then summed for the destination station \( j \). To make the index independent of the station’s overall demand, the sum of these weighted distances is normalized by the total number of trips received by station \( j \) (\( \sum_{i} T_{ij} \)). The resulting distance index provides a measure of how trip distances, adjusted for trip frequency, contribute to battery depletion at a station. In intuitive terms, higher \(DI_j\) means the station receives fewer trips overall and a greater share from distant origins, so bikes there spend more time docked recharging and less time in use. Because trips drain the battery while docking recharges it, this recharge effect dominates and higher \(DI_j\) is therefore associated with higher average battery levels. The formula of the index is as follows:

        \begin{equation}
          DI_j = \frac{\sum_{i} \frac{d_{i,j}}{n_i}}{\sum_{i} T_{ij}}
        \end{equation}

        The distance index was computed for all stations, and averaged over each neighborhood (Figure~\ref{correlations}C). As shown, neighborhoods in the northern areas tend to have a higher distance index, indicating that trips to these neighborhoods often originate from farther stations and that there are fewer incoming trips overall. This combination results in higher average battery levels. Conversely, neighborhoods in the central and southern parts of the city exhibit a lower distance index, suggesting that trips to these areas typically come from nearby stations and that there are more incoming trips, leading to lower battery levels. Zona Franca, as previously mentioned, deviates from this trend due to its unique industrial nature and low population density. The Pearson correlation between the distance index and average battery levels at the station level revealed a value of 0.47 (Figure~\ref{correlations}D), indicating a moderate positive correlation. This suggests that stations farther from others and receiving a lower number of trips tend to have higher average battery levels. Additionally, no significant collinearities were found between the inter-event time and the distance index (Appendix~\ref{app_paper2:subsection_collinearity}).


    \subsection[Modeling e-bikes' battery levels]{Modeling \acp{e-b}' battery levels}

        Using the battery model described in Section~\ref{sec:battery_modeling} (linear increase when docked, decrease during trips), the optimal parameters for the power consumption equation and charging speed were determined through hyperparameter optimization (see Section~\ref{sec:battery_optimization}). The fit between the observed and expected battery levels, using these optimal parameters, is shown in Figure \ref{fig:battery_fit}. In Figure \ref{fig:battery_fit}A, the predicted battery percentage is displayed, while in Figure \ref{fig:battery_fit}B, the predicted battery variation is presented. Despite the dispersion, a significant correlation is observed, with values closely aligned to the diagonal. Factors such as user heterogeneity, missing route information, or bikes not charging while docked could contribute to the dispersion. The parameters were also validated using the 8-hour dataset, with similar correlation values (Appendix~\ref{app_paper2:subsection_linear_battery_function}). In addition, a non-linear charging function was tested using a sigmoid-like function to account for faster charging times at lower battery levels. The fitting results are very similar to the linear function, with the optimal shape inferred being almost linear (Appendix~\ref{app_paper2:subsection_nonlinear_charging_function}).

        \begin{figure}[!htbp]
          \begin{center}
          \includegraphics[width=0.9\textwidth]{0_plots/paper2/main/battery_prediction.png}
          \end{center}
          \caption[Battery model validation using the 30-minute battery dataset]{
            \textbf{Battery model validation using the 30-minute battery dataset.} 
            \textbf{(A)} Observed battery level as a function of the predicted values. 
            \textbf{(B)} Battery variation observed as a function of the predicted values. Each point corresponds to one battery-level observation. The Root Mean Squared Error (RMSE) is $0.0178$ and the Pearson correlation coefficient is significant in both cases (p-val<0.001). The linear fit is shown in gray.
          } 
          \label{fig:battery_fit}
        \end{figure}

        This battery model can have two key applications. The first enables the assessment of battery consumption and destination probabilities for individual trips. The second integrates the Markov chain-based approach with the battery model to predict both \ac{e-b} locations and battery levels across the entire fleet.


    \subsection{Single trip application}

        Evaluating short-term scenarios for a single bike can be achieved by coupling the battery model with the distribution of possible destinations $P_{ij}^{w,h}$. In Figure \ref{fig:single_trip}, given \textit{St. Ciutat de Granada, 168} as the origin station, the probability of arriving at each destination is represented by the size of the dots, while the battery consumption for each bike trip is indicated by the color. The probability of destinations changes significantly between 8:00 and 18:00 despite the same origin station. At 18:00, trips tend to be shorter and head towards nearby stations, while morning trips are longer and converge towards a few work-related hot spots. Additionally, battery consumption varies notably based on the elevation changes along the route, with trips going downhill consuming less battery over the same distance. 

        \begin{figure}[!htbp]
          \begin{center}
          \includegraphics[width=0.9\textwidth]{0_plots/paper2/main/single_trip_daytype=0_station=42.pdf}
          \end{center}
          \caption[Single trip location and battery prediction]{
            \textbf{Single trip location and battery prediction.}
            Probability of finding a bike at a station and battery consumption of the trip when departing from the station in \textit{St. Ciutat de Granada, 168} on a weekday at \textbf{(A)} 08:00 and \textbf{(B)} 18:00. The origin station is highlighted with a triangle. The size of the dots corresponds to the probability of having that station as the destination and color corresponds to the battery consumed by the bike. The grey points correspond to destination stations with no trips.
          }
          \label{fig:single_trip}
        \end{figure}

        Furthermore, an impact index can be computed to assess the influence of each origin station on battery levels. This index depends on the probability of each destination given an origin station, the battery consumption of each possible trip, and the average hourly trips departing from station $i$ to other stations. The impact index can be written as:

        \begin{equation}
        \gamma^{w,h}_i=\langle T^{w,h}_i \rangle P^{w,h}_{ij} \Delta B_{ij}
        \end{equation} 

        where $\langle T^{w,h}_i \rangle$ represents the average number of trips departing from station $i$, $P^{w,h}_{ij}$ is the probability that a trip generated in station $i$ ends in station $j$, and $\Delta B_{ij}$ is the battery consumed by a trip from $i$ to $j$. An example of the impact index $\gamma$ at 8:00 and 18:00 on a regular weekday is shown in Figure \ref{fig:impact}. Due to the variation in trip patterns captured by $P^{w,h}_{ij}$, the stations with the highest impact change throughout the day.

        \begin{figure}[!htbp]
          \begin{center}
          \includegraphics[width=0.9\textwidth]{0_plots/paper2/main/impact_index_0.pdf}
          \end{center}
          \caption[Station impact index at different hours of the day]{
            \textbf{Station impact index at different hours of the day.} Impact index of the stations at \textbf{(A)} 08:00 and \textbf{(B)} 18:00. The station with higher impact at each hour is highlighted with a larger dot size. In the morning, the \textit{St. Provença, 125} station has the highest impact, while in the afternoon, \textit{St. Ciutat de Granada, 168} has the greatest impact.
            } 
          \label{fig:impact}
        \end{figure}

    \subsection{Coupling the Markov chain and the battery models}

        The adjusted battery consumption function can be integrated into the Markov chain approach to provide joint predictions of both location and battery levels over multiple trips, either at the single-bike or \ac{bss} level. Inter-event times are used to calculate the charging percentage, while the distance and duration of trips are used to calculate energy consumption. Thus, the inputs required for the coupled model are the location and battery level of each bike along with the day of the week and the hour, which together define the starting point of the simulation.

        To validate this approach, systematic simulations were performed using initial conditions from the 2023 mobility data and 30-minute battery data (Figure \ref{fig:markov_correlations}). Due to the inherent noise in micro-mobility service usage (i.e., the unpredictability of which bike at a station will be selected by the next user), the validation was conducted at the station level. For each validation snapshot, the initial conditions were determined by the number of available bikes and their average battery level at each station in the snapshot. To address potential gaps in battery data, the battery level of each bike was set to the station average. However, while comparisons between actual and predicted values were made at the station level, the mobility and battery dynamics of each bikes were simulated individually. Additionally, to prevent data leakage, the probabilities of the Markov model were based on data from the year 2022.

        \begin{figure}[!htbp]
          \begin{center}
          \includegraphics[width=0.9\textwidth]{0_plots/paper2/main/markov_correlations.pdf}
          \end{center}
          \caption[Correlations for the Markov chain model across prediction horizons.]{
            \textbf{Correlations for the Markov chain model across prediction horizons.} 
            \textbf{(A)} Pearson correlation coefficient between actual and predicted bike availability per station as a function of $\Delta t$. 
            \textbf{(B)} Pearson correlation coefficient between actual and predicted average battery level per station as a function of $\Delta t$.
          }
          \label{fig:markov_correlations}
        \end{figure}

        Once both the mobility and battery consumption models were validated, simulations across various initial scenarios and time horizons were conducted to demonstrate the potential applications of our framework (Appendix~\ref{app_paper2:section_initial_conditions_for_battery_predictions}). For these predictions, the initial conditions were derived from the average number of available bikes and their battery levels across stations in snapshots corresponding to the same type of day and hour. First, it is possible to compute the probability of finding a bike at a given destination and its battery level as a function of its origin, following one or more trips within a time period $\Delta t$. The final battery levels are determined by averaging over all the Markov iterations. As an example, Figure \ref{fig:markov_trips} shows the probability and battery levels of bikes departing from the station at Glòries (\textit{St. Ciutat de Granada, 168}) at 08:00 after \textbf{(A)} 30 and \textbf{(B)} 120 minutes. The role of distance from the origin is evident in the short-term predictions, with lower battery levels at greater distances and inclinations. However, this effect diminishes over longer time windows, as multiple trips can occur before reaching a station. For longer periods, a greater variety of final destinations is also possible.


        \begin{figure}[!htbp]
          \begin{center}
          \includegraphics[width=0.9\textwidth]{0_plots/paper2/main/markov_trips_daytype=0-hour=08-station=42.pdf}
          \end{center}
          \caption[Predicted bike locations and battery levels for trips starting at 8:00 from Glòries station]{
            \textbf{Predicted bike locations and battery levels for trips starting at 8:00 from Glòries station.} Probability of finding a bike and average battery levels after \textbf{(A)} 30 minutes and \textbf{(B)} 120 minutes of simulation for trips departing from the Glòries station (St.\ Ciutat de Granada, 168) at 08:00 on an average weekday. The origin station is highlighted with a triangle.
          }
          \label{fig:markov_trips}
        \end{figure}

        As a second application, the Markov chain approach allows to transition from bike-level to system-level predictions. To calculate the number of available bikes per station across the entire fleet, the probabilities $f^{w,h}_{i,j}(\Delta t)$ for each destination station $j$ are summed, which may result in non-integer values. Battery levels, on the other hand, are calculated by multiplying each bike's probability by its corresponding battery level and then dividing by the number of available bikes. Figure~\ref{fig:markov_battery} displays the number of available bikes and the average battery level after two hours of simulation for two different starting times, 08:00 and 18:00. Stations near the city center exhibit lower battery levels, consistent with the data analysis results. Meanwhile, the northern part of the city, which experiences lower usage rates, shows higher battery levels. Additionally, lower battery levels are observed in the afternoon due to increased battery consumption throughout the day.

        \begin{figure}[!htbp]
          \begin{center}
          \includegraphics[width=0.9\textwidth]{0_plots/paper2/main/markov_battery_daytype=0-dt=120.pdf}
          \end{center}
          \caption[System-level predictions of bike availability and battery levels after 120 minutes of simulation]{
            \textbf{System-level predictions of bike availability and battery levels after 120 minutes of simulation.} 
            Predictions combine the Markov mobility model with the battery equations for simulations starting at \textbf{(A)} 08:00 and \textbf{(B)} 18:00 on an average weekday.
          }
          \label{fig:markov_battery}
        \end{figure}

\section{Discussion}
\label{sec:disc_battery}

This Section reflects on the contribution on modeling fleet mobility and \ac{e-b} battery status developed in Part~\ref{pt:predictive_operations}. It builds on the modeling setup and empirical analysis presented in Chapter~\ref{ch:mobility_and_battery_modeling_bicing}, and discusses the main implications of the proposed framework for battery-aware operations in docked \acp{bss}. Methodological details and results are provided in Sections~\ref{sec:battery_methodology} and~\ref{sec:battery_results}, respectively.

A central finding of this contribution is the presence of pronounced spatiotemporal variation in station-level battery levels. Battery levels were generally lower in the evenings, reflecting the impact of daily usage, and higher in the mornings after overnight resting periods. Spatially, stations in northern neighborhoods exhibited higher average battery levels, likely attributed to lower usage rates and longer resting times. Additionally, notable correlations were observed between stations' battery levels, their inter-event times, and their distance from other stations.

Building on these observations, the developed distance index provides valuable insights into one of the primary operational challenges faced by electric \acp{bss}: battery level management. Stations with lower distance index values are more prone to accelerated battery depletion, as they tend to receive more frequent trips from nearby origins, leaving less time for recharging between trips. From an operational standpoint, this metric offers \ac{bss} operators a practical tool to identify such stations and prioritize interventions, including more frequent battery replacements, recharging, or strategic bike redistribution, thereby enhancing the system's efficiency and reliability.

Methodologically, the presented bike battery forecasting approach combines a traditional probabilistic framework with heterogeneous data sources, resulting in a multifaceted tool capable of improving the real-time management of micro-mobility services. By coupling a validated battery dynamics function with a Markov-chain mobility model, the framework enables joint predictions about future \ac{e-b} availability and battery status. This is particularly relevant in increasingly electrified fleets, where predicting availability alone can be insufficient: a docked bike with a depleted battery is effectively unavailable. The main value of the approach is therefore its ability to move toward microscopic forecasts that capture not only where bikes are likely to be, but also whether they will be usable upon arrival.

At the same time, the scope and accuracy of the results are conditioned by several data limitations. The available battery data were limited in both resolution and coverage: one dataset had an 8-hour temporal resolution, which complicates robust model training and evaluation, while the 30-minute resolution dataset covered only two weeks and did not include bikes in transit. These constraints reduce the ability to learn fine-grained dynamics and to fully validate predictions under a wide range of operating conditions.

In addition, the lack of detailed trip information restricts the ability to model battery consumption more realistically. Because the exact routes taken by users were unavailable, battery depletion could not be estimated from realized paths, and round trips could not be reconstructed. Further sources of heterogeneity that affect consumption are also unobserved, including bicycle condition (e.g., tire inflation) and rider characteristics (e.g., weight or height). These omissions are common in operational datasets, but they imply that the battery consumption component must rely on simplifying assumptions and calibrated parameters rather than on fully observed on-route behavior.

These limitations also point to clear directions for improving the approach. Access to \ac{gps}-based route information would likely improve power-consumption estimates by allowing the battery function to be applied along realized trajectories, while higher-resolution and longer time series could enhance the battery predictions. Additionally, experimental data collected through ad-hoc studies could improve the model's performance by fine-tuning the parameters, ensuring a more accurate representation of real-world conditions.

Overall, this contribution shows that battery-aware forecasting can support operations by anticipating not only where bikes will be, but also whether they will be usable. Despite the limited resolution and missing route-level information, the results demonstrate that meaningful decision support is feasible with already available \ac{bss} data and would further improve with richer trip observations.

